\documentclass[11pt,a4paper]{article}
\usepackage{fontspec}
\setmainfont{DejaVu Sans}
\setsansfont{DejaVu Sans}
\usepackage{polyglossia}
\setdefaultlanguage{russian}
\setotherlanguage{english}
\usepackage{amsmath,amssymb,mathtools}
\usepackage{geometry}
\usepackage{xcolor}
\usepackage{enumitem}
\usepackage{microtype}
\usepackage{fancyhdr}
\usepackage{tcolorbox}
\usepackage{hyperref}
\geometry{left=18mm,right=18mm,top=18mm,bottom=20mm}
\setlength{\parindent}{0pt}
\setlength{\parskip}{4pt}
\definecolor{hseblue}{HTML}{003D7C}
\definecolor{softblue}{HTML}{EEF4FA}
\hypersetup{colorlinks=true,linkcolor=hseblue,urlcolor=hseblue}
\pagestyle{fancy}
\fancyhf{}
\lhead{Методы оптимизации в машинном обучении}
\rhead{Семинар 2}
\cfoot{\thepage}
\newcommand{\R}{\mathbb{R}}
\newcommand{\norm}[1]{\left\lVert #1\right\rVert}
\newcommand{\grad}{\nabla}
\newcounter{problem}
\newenvironment{problem}[1][]{%
  \refstepcounter{problem}%
  \begin{tcolorbox}[colback=white,colframe=hseblue!70!black,boxrule=0.7pt,arc=2mm,left=2.5mm,right=2.5mm,top=2mm,bottom=2mm,title={\textbf{Задача \theproblem.} #1},fonttitle=\normalsize]
}{\end{tcolorbox}}
\begin{document}

{\LARGE\bfseries\color{hseblue} Семинар 2. Геометрия оптимизации}

\vspace{0.4em}
На эту часть семинара отводится примерно 40 минут. В задачах важны не только вычисления, но и короткая интерпретация результата.

\begin{problem}[Скалярное произведение и геометрия направлений]
Даны векторы
\[
x=(1,2,-1),\qquad
v_1=(2,0,1),\qquad
v_2=(1,-1,-1),\qquad
v_3=(-2,-1,0).
\]
\begin{enumerate}[label=\alph*)]
\item Для каждого $v_i$ вычислите $x^\top v_i$.
\item Определите, образует ли $v_i$ с $x$ острый, прямой или тупой угол.
\item Придумайте свой ненулевой вектор $u$, ортогональный $x$, и проверьте это вычислением.
\item Сформулируйте, что знак скалярного произведения говорит о взаимном направлении двух векторов.
\end{enumerate}
\end{problem}

\begin{problem}[Локальная линейная модель]
Рассмотрим функцию
\[
f(x,y)=x^2+xy+y^2
\]
и точку $x_0=(1,1)$.
\begin{enumerate}[label=\alph*)]
\item Найдите $\grad f(x_0)$ и $f(x_0)$.
\item Используя локальную линейную модель
\[
f(x_0+h)\approx f(x_0)+\grad f(x_0)^\top h,
\]
приблизьте значение $f(1.01,0.98)$.
\item Вычислите точное значение $f(1.01,0.98)$ и абсолютную ошибку приближения.
\item Почему разумно ожидать, что при уменьшении длины $h$ линейное приближение становится точнее?
\end{enumerate}
\end{problem}

\begin{problem}[Направленные производные]
Пусть
\[
f(x,y)=x^2+4y^2,\qquad x_0=(1,1),
\]
и заданы единичные направления
\[
v_1=(1,0),\qquad
v_2=(0,1),\qquad
v_3=\frac{1}{\sqrt2}(1,-1).
\]
\begin{enumerate}[label=\alph*)]
\item Найдите $\grad f(x_0)$.
\item Вычислите $D_{v_i}f(x_0)$ для каждого направления.
\item В каких из этих направлений функция локально возрастает, а в каких убывает?
\item Какое из трёх направлений даёт наиболее быстрый локальный рост? Наиболее быстрое локальное убывание?
\end{enumerate}
\end{problem}

\begin{problem}[Направление наискорейшего роста]
Пусть $f:\R^d\to\R$ дифференцируема в точке $x$ и $\grad f(x)\neq0$. Для любого единичного направления $v$
\[
D_vf(x)=\grad f(x)^\top v.
\]
\begin{enumerate}[label=\alph*)]
\item Используя неравенство Коши--Буняковского, докажите, что
\[
D_vf(x)\le \norm{\grad f(x)}_2.
\]
\item Покажите, что равенство достигается при
\[
v=\frac{\grad f(x)}{\norm{\grad f(x)}_2}.
\]
\item Сформулируйте и докажите аналогичное утверждение о направлении наискорейшего локального убывания.
\end{enumerate}
\end{problem}

\begin{problem}[Линии уровня и градиент]
Рассмотрим
\[
f(x,y)=x^2+4y^2.
\]
Линия уровня $f(x,y)=5$ проходит через точку $x_0=(1,1)$.
\begin{enumerate}[label=\alph*)]
\item Найдите $\grad f(x_0)$.
\item Найдите ненулевой вектор $v$, касательный к линии уровня в точке $x_0$.
\item Проверьте, что $\grad f(x_0)^\top v=0$.
\item Запишите уравнение касательной прямой к эллипсу $x^2+4y^2=5$ в точке $(1,1)$.
\item Объясните геометрический смысл полученной ортогональности.
\end{enumerate}
\end{problem}

\begin{problem}[Кривизна вдоль направления]
Пусть в некоторой точке гессиан функции равен
\[
H=
\begin{pmatrix}
2&0\\
0&8
\end{pmatrix}.
\]
Рассмотрите единичные направления
\[
v_1=(1,0),\qquad
v_2=(0,1),\qquad
v_3=\frac{1}{\sqrt2}(1,1).
\]
\begin{enumerate}[label=\alph*)]
\item Вычислите направленную кривизну $\kappa(v)=v^\top Hv$ для каждого $v_i$.
\item В каком направлении поверхность наиболее «крутая» из трёх? В каком наиболее «плоская»?
\item Как связаны найденные значения с собственными значениями матрицы $H$?
\end{enumerate}
\end{problem}

\begin{problem}[Собственные значения и квадратичная форма]
Пусть
\[
A=
\begin{pmatrix}
3&1\\
1&3
\end{pmatrix}.
\]
\begin{enumerate}[label=\alph*)]
\item Найдите собственные значения $A$ и соответствующие собственные направления.
\item Определите знак квадратичной формы $x^\top Ax$ для $x\neq0$.
\item Докажите для этой матрицы оценку
\[
2\norm{x}_2^2\le x^\top Ax\le4\norm{x}_2^2.
\]
\item Как выглядит поверхность $f(x)=\frac12x^\top Ax$ около начала координат: как «чаша», «купол» или «седло»? Почему?
\end{enumerate}
\end{problem}

\begin{problem}[Стационарные точки и их классификация]
Рассмотрим функцию
\[
f(x,y)=x^2+xy+y^2-2x.
\]
\begin{enumerate}[label=\alph*)]
\item Найдите все стационарные точки из условия $\grad f(x,y)=0$.
\item Вычислите гессиан и его собственные значения.
\item Классифицируйте найденную стационарную точку.
\item Теперь рассмотрите
\[
g(x,y)=x^2-y^2.
\]
Покажите, что $(0,0)$ является стационарной точкой, и классифицируйте её по гессиану.
\end{enumerate}
\end{problem}

\end{document}
